\begin{table}[htbp]
\centering
\scriptsize
\caption{Exploratory diagnosis of the non-finite KalmanNet loss (one training batch of the default configuration, untrained network). Source: \path{experiments/filter_family/results/lorenz_knet_nan_diagnosis.csv}.}\label{tab:lorenz_diag}
\resizebox{\linewidth}{!}{%
\begin{tabular}{lp{0.55\linewidth}}
\toprule
Quantity & Value \\
\midrule
training data finite & True \\
first step with a non-finite value (0-based) & 73 \\
non-finite quantities at that step & m1x\_prior;m1y;KGain;m1x\_posterior \\
affected batch rows & 24 \\
first sub-modules with non-finite output & FC7;GRU\_S;FC2;FC3;FC4 \\
max abs state of the affected row at the previous step & $7.658\times10^{17}$ \\
float32 maximum & $3.403\times10^{38}$ \\
first non-finite Taylor term (order j) & 5 \\
operation & torch.matrix\_power(A*delta\_t, j) in f (Simulations/Lorenz\_Atractor/parameters.py:90) \\
batch loss finite (anomaly run) & False \\
anomaly detection message & Function 'MseLossBackward0' returned nan values in its 0th output. \\
PyTorch version & 2.15.0+cpu \\
\bottomrule
\end{tabular}}
\end{table}
